% القسم: المناهج
% الفصل: البراهين
% المبحث: لا أستطيع

\documentclass[../../../include/open-logic-section]{subfiles}

\begin{document}

\olfileid{mth}{prf}{cnt}

\olsection{إذا عسر البرهان}

لا يخلو طالب من ساعة يدعوه فيها العسر إلى ترك المسألة؛ غير أنك
\emph{تستطيع} أن تتمها. وعندك مدرسك ومساعده وزملاؤك، وكلهم قادر على
أن يعينك، فلا تمتنع من سؤالهم. وهذه وصايا تدفع بلوغ الأمر حد الأزمة،
وتدل على ما يصنعه من همّ بالاستسلام.

ومن أسباب النجاح في حل المسائل ما يأتي:
\begin{enumerate}
\item \emph{بادر ما استطعت.} أشغال الفصل كثيرة، والنفس تميل إلى
  التسويف؛ فخير ما تستعين به أن تشرع في الواجب قبل ضيق الوقت. فإن
  وقفت عند موضع وجدْت فسحة للبحث عن مخرجه، ولم يبق حظك مجرد الأسف.
\item \emph{ذاكر زملاءك.} فلست منفردًا بهذا العناء؛ وقد تعرض لغيرك
  صعوبة، وإن كانت في موضع غير موضعك. ومذاكرة القرين تريك المسألة من
  وجه آخر فينكشف به المغلق. ولا تنقل حلّه، بل استرشده بتلميح، أو صف له
  موضع توقفك واسأله عن الخطوة التالية. فإذا فهمت فأعن غيرك كما أُعنت؛
  فإن شرح الفكرة للآخر يزيدك بها فهمًا.
\item \emph{استعن بمن يعلم.} بين يديك موارد كثيرة، ومدرسك ومساعده
  إنما هما هناك لعونك، وهما \emph{يريدان} نجاحك. وبوسعهما أن يرشداك
  في حل المسألة، وأن يعيّنا المرحلة التي استغلقت عليك.
\item \emph{استرح قليلًا.} فقد يكون طول مكثك على المسألة هو الذي
  حجب وجهها. اقطع النظر عنها زمنًا قصيرًا، واشرب قدحًا من الشاي، أو
  انتقل إلى مسألة أخرى، ثم عد إليها بذهن مستريح؛ وربما نفع أن تبيت
  عليها ليلة.
\end{enumerate}

وهذه الوصايا كلها مبنية على المبادرة؛ فمن لم يبدأ البرهان إلا في
الثانية من صباح اليوم السابق لموعد تسليمه ضاقت عليه أكثر أبواب النفع.

وليس هذا دعوة إلى التشاؤم؛ فإن عناء البرهان يعقبه ثمرته. وقد اتخذه
قوم صناعة عمرهم، فلا بد أن فيه لذة تدعوهم إليه. وحل المسائل وإقامة
البراهين مهارة كسائر المهارات، لا تنال إلا بكثرة المران. فإذا رأيت
زميلًا يسهل عليه ما عسر عليك، فلعله سبقك إلى الممارسة؛ فلا تعجل إلى
ترك العمل.

فإن نفد وقتك أو صبرك في مسألة بعينها فلا حرج. وليس ذلك دليل بلادة،
ولا حكمًا بأن فهمها ممتنع عليك أبدًا. سل مدرسك أو زميلك عن طريق الحل،
ثم عيّن موضع الخطأ أو الوقوف كي تتجنبه إذا عاد نظيره. وبعد أن تفهم
الحل حاول إعادته من غير نظر إليه. وفي المرة الآتية بادر، واطلب العون،
قبل أن يضيق الوقت.

\end{document}
